\section{Introducción}

Se pretende desarrollar un sistema similar a Airbnb donde personas particulares pueden reservar alojamientos de propietarios particulares. Tanto los inquilinos como propietarios deben registrarse previamente para utilizar el sistema.
\begin{itemize}
    \item Los propietarios (una vez logeados) pueden dar de alta propiedades (viviendas completas o habitaciones individuales) ofertándolas para ser alquiladas. Los propietarios marcan si desean que los usuarios realicen reservas directas, o guardarse el derecho de confirmar o no las solicitudes de reserva. 

    \item Los usuarios pueden primero buscar alojamiento (destino, fechas, tipo de inmueble, etc.) para examinar las opciones de alquiler. Excepcionalmente también pueden aplicar filtros avanzados para, por ejemplo: ver aquellos inmuebles con posibilidad de reserva inmediata; seleccionar filtros para ciertas comodidades; o seleccionar política de cancelación de reserva.

    \item Para realizar la búsqueda de inmuebles no es necesario estar logeado, pero los usuarios registrados pueden agregar algunas propiedades a su lista de deseo.

    \item Los inquilinos logeados, una vez han seleccionado el inmueble a alquilar, pueden completar la reserva. Si la propiedad permite reserva inmediata se puede completar el pago directamente (vía tarjeta de crédito, débito o paypal).

    \item Si la propiedad o inmueble no permite reserva inmediata se hace una solicitud de reserva (completando también en este momento el pago). La solicitud de reserva es notificada al propietario quién puede confirmar o no la reserva. Si la confirmación es positiva, esto se le notifica al inquilino. Si la solicitud es denegada, el sistema devuelve el dinero al inquilino y es igualmente notificado al cliente. 
\end{itemize}

En este documento se describirá todo lo necesario para poder desarrollar la aplicación web llamada Bckroom. El documento contendrá las historias de usuario, arquitectura y diseños de la aplicación, un control de configuración y control de la calidad.

Para este proyecto se realizará con la metología ágil SCRUM. Los sprints tendrán una duración de 3 semanas, las reuniones semanales se realizarán los viernes y las reuniones daily tendrán una duración máxima de 15 minutos. También se realizarán reuniones al finalizar cada sprint de una duración máxima de 2 horas y 30 minutos y mínima de 30 minutos en la que se analizará las funcionalidades implementadas, los problemas surgidos con las funcionalidades y comentar el retraso de las funcionalidades, si se producen.